% Part: first-order-logic
% Chapter: proof-systems
% Section: natural-deduction

\documentclass[../../../include/open-logic-section]{subfiles}

\begin{document}

\iftag{FOL}
      {\olfileid{fol}{prf}{ntd}}
      {\olfileid{pl}{prf}{ntd}}

\olsection{Deduksi Natural}

Deduksi natural iku sistem !!a{derivation} kang dirancang kanggo
nyerupani penalaran nyata (mligine jinis penalaran tertata kang
digunakake matematikawan). Penalaran nyata lumaku nganggo sawetara
pola ``natural''. Tuladhane, bukti miturut kasus ngidini kita netepake
kesimpulan adhedhasar premis disjungtif, yaiku kanthi netepake yen
kesimpulan ndherek saka saben disjung. Bukti ora langsung ngidini kita
netepake kesimpulan kanthi nuduhake yen negasine nuwuhake kontradiksi.
Bukti kondisional netepake pratelan kondisional ``yen \dots mula
\dots'' kanthi nuduhake yen konsekuen ndherek anteseden. Deduksi
natural iku formalisasi sawetara inferensi natural iki. Saben konektif
logis lan kuantor duwe rong aturan, yaiku aturan introduksi lan
eliminasi, lan saben aturan cocog karo sawijining pola inferensi
natural. Tuladhane, $\Intro{\lif}$ cocog karo bukti kondisional, lan
$\Elim{\lor}$ karo bukti miturut kasus. Salah siji aturan kang mligi
prasaja yaiku $\Elim{\land}$, kang ngidini inferensi saka $!A \land !B$
menyang~$!A$ (utawa $!B$).

Salah siji ciri kang mbedakake deduksi natural saka sistem
!!{derivation} liyane yaiku panganggone asumsi. !!^a{derivation} ing
deduksi natural iku wit !!{formula}s. Ana siji !!{formula} ing oyod wit
!!{formula}s, lan ``godhong'' wit iku !!{formula}s kang dadi dhasar
derivasi kesimpulan. Ing deduksi natural, sawetara !!{formula}s godhong
duwe peran ing njero !!{derivation}, nanging wis ``entek digunakake''
nalika !!{derivation} tekan kesimpulan. Iki cocog karo laku ing
penalaran nyata, yaiku ngenalake hipotesis kang mung lumaku sedhela.
Tuladhane, ing bukti miturut kasus, kita nganggep saben disjung bener;
ing bukti kondisional, kita nganggep anteseden bener; ing bukti ora
langsung, kita nganggep negasi saka kesimpulan bener. Cara ngenalake
asumsi hipotetis banjur nyingkirake asumsi iku kanggo netepake tahap
antara dadi ciri khas deduksi natural. Formula ing godhong
!!{derivation} deduksi natural diarani asumsi, lan sawetara aturan
inferensi bisa ``!!{discharge}'' asumsi kasebut. Tuladhane, yen kita
duwe !!a{derivation} saka~$!B$ adhedhasar sawetara asumsi kang
ngemot~$!A$, aturan $\Intro{\lif}$ ngidini kita nyimpulake~$!A \lif !B$
lan mbubarake saben asumsi kanthi wujud~$!A$. (Kanggo nyathet asumsi
endi kang dibubarake ing inferensi endi, kita menehi angka minangka
label kanggo inferensi lan asumsi kang dibubarake.) Asumsi kang isih
!!{undischarged} ing pungkasan !!{derivation} bebarengan cukup kanggo
njamin kabeneran kesimpulan, mula !!a{derivation} netepake yen asumsi
!!{undischarged} kasebut ngakibatake kesimpulane kanthi semantis.

Relasi $\Gamma \Proves !A$ adhedhasar deduksi natural lumaku yen lan
mung yen ana !!a{derivation} kang nduweni $!A$~minangka !!{sentence}
pungkasan ing wit, lan saben godhong kang !!{undischarged} kalebu
ing~$\Gamma$. $!A$~iku teorema ing deduksi natural yen lan mung yen ana
!!a{derivation} kang nduweni $!A$~minangka !!{sentence} pungkasan lan
kabeh asumsi wis !!{discharged}. Tuladhane, iki !!a{derivation} kang
nuduhake yen $\Proves (!A \land !B) \lif !A$:
\begin{prooftree}
  \AxiomC{$\Discharge{!A \land !B}{1}$}
  \RightLabel{\Elim{\land}}
  \UnaryInfC{$!A$}
  \DischargeRule{\Intro{\lif}}{1}
  \UnaryInfC{$(!A \land !B) \lif !A$}
\end{prooftree}
Label~$1$ nuduhake yen asumsi $!A \land !B$ wis !!{discharged} ing
inferensi \Intro{\lif}.
  
Himpunan~$\Gamma$ ora konsisten yen lan mung yen
$\Gamma \Proves \lfalse$ ing deduksi natural. Aturan \FalseInt{}
ndadekake saben !!{sentence} bisa !!{derive}d saka himpunan kang ora
konsisten.

Sistem deduksi natural dikembangake dening Gerhard Gentzen lan
Stanis\l{}aw Ja\'skowski ing taun 1930-an, banjur dikembangake maneh
dening Dag Prawitz lan Frederic Fitch. Amarga inferensine nyerupani
metodhe bukti natural, sistem iki disenengi para filsuf. Versi kang
dikembangake Fitch kerep digunakake ing buku teks logika pambuka. Ing
filsafat logika, aturan deduksi natural kadhangkala dianggep menehi
makna operator logis (``semantik teoretis-bukti'').

\end{document}
